%%% SECTION 3: CALIFORNIA AND FEDERAL REGULATORY LANDSCAPE %%%

California serves as the site of the first Physical AI oncology clinical trial,
making a dual-jurisdiction analysis of state and federal law essential for
understanding the complete regulatory environment. This section adapts the
comprehensive California and federal legal framework analysis to Physical AI
oncology trial contexts, providing the regulatory map that trial sponsors,
investigators, and sites must navigate.

\subsection{Regulatory Classification Framework}
\label{subsec:regulatory-classification}

The regulatory landscape for Physical AI in oncology trials can be organized
along two classification axes. The first axis distinguishes between AI systems
that qualify as regulated medical products and those that function as research
operational tools without direct medical product classification. The second axis,
applicable to device-classified Physical AI systems, distinguishes between
investigational and approved status, and between significant risk and
non-significant risk designations.

The federal device definition under 21 U.S.C. \S~321(h) encompasses instruments,
machines, and other articles intended for use in diagnosis, treatment, or
prevention of disease, which includes many Physical AI system components. The
investigational device exemption authority under 21 U.S.C. \S~360j(g)
establishes procedures for device use during clinical investigation. The FD\&C
Act \S~520(o) software exclusions from the 21st Century Cures Act exclude
certain software functions from the device definition, but these exclusions were
designed primarily for software-only clinical decision support and may not
encompass the full range of Physical AI functionalities that involve direct
physical patient interaction.

Physical AI systems present classification challenges because a single system
may simultaneously occupy multiple categories. A surgical robot performing a
tumor resection functions as a medical device. The AI algorithm guiding its
surgical decisions may qualify as software as a medical device (SaMD). The data
collection and transmission functions may be classified as electronic data
systems under Part 11. The robot's interaction with the patient triggers human
subjects protection requirements. This multi-category nature underscores the
need for the integrated regulatory approach provided by this National Platform.

\subsection{California Regulatory Authority}
\label{subsec:california-authority}

\subsubsection{Human Subjects Protections}

California's Protection of Human Subjects in Medical Experimentation Act (Cal.
Health \& Safety Code \S\S~24170-24179.5) establishes protections for research
participants that supplement federal requirements. The Experimental Subject's
Bill of Rights (\S~24172) enumerates specific rights that must be communicated
to research participants, creating obligations beyond federal informed consent
requirements. The definition of ``medical experiment'' (\S~24174) is broad
enough to encompass Physical AI procedures where robots perform functions that
constitute medical experimentation.

For Physical AI trial participants in California, these provisions create
additional protections. The Experimental Subject's Bill of Rights requires
disclosure of the nature of the experiment (including robotic involvement),
risks and discomforts (including robot malfunction scenarios), and the right to
withdraw. The Clinical Trial Site Complete Documentation \cite{site-docs}
incorporates these California-specific requirements into its documentation
package, and AB 2847 from that package establishes additional patient rights
specific to Physical AI interactions.

\subsubsection{Health Data Privacy and Security}

The Confidentiality of Medical Information Act (CMIA, Cal. Civ. Code
\S\S~56-56.37) governs the handling of medical information by healthcare
providers and other entities. Electronic medical information safeguards
(\S~56.101) require specific protections for electronically maintained medical
data, directly applicable to the digital records generated by Physical AI
systems. California's data breach notification requirements (Cal. Civ. Code
\S~1798.82) and medical information breach requirements (Cal. Health \& Safety
Code \S~1280.15) impose notification obligations when Physical AI system
security breaches compromise patient data.

Physical AI systems generate new categories of health data including robot
telemetry, sensor readings, video recordings of procedures, and AI decision
logs. CMIA requirements apply to these data categories when they constitute
medical information, creating compliance obligations that interact with federal
HIPAA \cite{hipaa} protections. The Physical AI Federated Learning framework
\cite{fl-paper} addresses these dual privacy requirements through
privacy-preserving approaches that minimize data centralization.

\subsubsection{Consumer Privacy}

CCPA/CPRA (Cal. Civ. Code \S\S~1798.100 et seq.) provides broad consumer
privacy protections with specific exemptions for clinical trial and biomedical
research data. Clinical trial data may be statutorily exempt in many common
trial contexts, but non-exempt data streams from Physical AI systems may still
trigger CCPA obligations. For example, if a Physical AI system collects data
about a patient's movements or preferences that falls outside the clinical trial
data exemption, CCPA rights to access, deletion, and opt-out would apply.

\subsubsection{AI-Specific Healthcare Communications}

AB 3030 (Cal. Health \& Safety Code \S~1339.75) requires disclosure when
AI-generated content is used in patient clinical communications. AB 489
addresses health advice from artificial intelligence. Physical AI systems that
communicate with patients, whether through social companion robots providing
emotional support or through AI-generated treatment explanations, must comply
with these disclosure requirements. The patient instructions documented in
Patient Instructions: Physical AI Trials \cite{patient-instructions-paper}
incorporate appropriate disclosures for each of the ten robot types.

\subsubsection{Medical Practice and Licensing}

California's unlicensed practice of medicine provisions (Cal. Bus. \& Prof.
Code \S~2052) raise questions about Physical AI systems that perform functions
traditionally reserved for licensed practitioners. When a surgical robot
performs tissue manipulation or a needle-placement robot conducts a biopsy, the
boundary between machine-assisted medical practice and autonomous machine
practice becomes legally significant. The adapted 21 CFR Part 312
\cite{cfr312-adapt} addresses this through its Physical AI operator definition,
which establishes that qualified human oversight must be maintained even when
robots perform physical procedures.

\subsection{Federal Regulatory Authority}
\label{subsec:federal-authority}

\subsubsection{Investigational Product Frameworks}

The IND regulations (21 CFR Part 312) \cite{cfr-part-312} govern investigational
drug use in clinical trials. When Physical AI systems are used alongside
investigational drugs, the IND framework must accommodate the additional
complexity of robotic drug administration, AI-guided dosing, and automated
safety monitoring. The Physical AI Adaption: 21 CFR Part 312 \cite{cfr312-adapt}
provides the adapted framework spanning ten subparts, including the Phase 0
simulation validation stage that requires demonstration of Physical AI
system safety before human subjects enrollment.

IDE regulations (21 CFR Part 812) govern investigational device use. When
Physical AI systems themselves qualify as investigational devices, sponsors must
obtain IDE approval or qualify for exemption. The interaction between IND and
IDE requirements becomes particularly complex when Physical AI systems are used
to administer investigational drugs, potentially triggering both regulatory
pathways simultaneously.

\subsubsection{Human Subjects Protection}

Federal human subjects protection operates through 21 CFR Part 50 \cite{cfr-part-50}
(FDA-specific protections), 21 CFR Part 56 (IRB requirements), and the Common
Rule (45 CFR Part 46). Physical AI adds new dimensions to informed consent
because participants must understand their interactions with robotic systems, AI
decision-making processes, safety monitoring systems, and their right to refuse
or pause robotic procedures. The Physical AI Adaption: 21 CFR Part 50
\cite{cfr50-adapt} addresses these dimensions through 17 Physical AI-specific
definitions and a new Subpart C dedicated to Physical AI safety requirements.

The NCI Central Institutional Review Board \cite{nci-cirb} provides a model for
centralized IRB review that could be adapted for Physical AI trial oversight,
enabling consistent evaluation of Physical AI risks and protections across
multiple trial sites.

\subsubsection{Electronic Records and Data Integrity}

21 CFR Part 11 governs electronic records and electronic signatures in
FDA-regulated activities. Physical AI systems generate, store, and transmit
electronic clinical trial records including robot telemetry, sensor data,
autonomous decision logs, imaging data, and treatment records. These constitute
new categories of electronic records that must comply with Part 11 requirements
for trustworthiness and reliability.

Interoperability standards including FHIR \cite{hl7-fhir} for healthcare data
exchange and DICOM \cite{dicom} for medical imaging provide standardized
formats for some Physical AI data streams. The MCP server infrastructure
\cite{national-mcp-paper} extends these standards by providing standardized
communication protocols for Physical AI-specific data that existing standards
do not cover, including real-time robot telemetry, task order management, and
safety interlock states.

\subsubsection{Privacy and Security}

HIPAA Privacy and Security Rules (45 CFR Parts 160 and 164) \cite{hipaa} govern
the use and disclosure of protected health information (PHI) by covered
entities and business associates. Physical AI systems handling PHI must comply
with HIPAA requirements, creating obligations for data encryption, access
controls, audit trails, and breach notification. The federated learning
framework \cite{fl-paper} and federated learning in healthcare
\cite{federated-learning-healthcare} provide privacy-preserving approaches
that reduce the need to centralize PHI while enabling multi-site model
improvement.

\subsection{Robotics-Specific Regulatory Considerations}
\label{subsec:robotics-regulatory}

FDA's position on robotically-assisted surgical (RAS) systems in oncology
presents a specific regulatory consideration. RAS systems have not received
marketing authorization specifically for cancer prevention or treatment,
although they are used in oncology-related procedures under existing
clearances. The significance of oncology-related endpoints (recurrence,
disease-free survival, overall survival) for Physical AI trials means that
trial designs must demonstrate not just procedural safety but also oncology
outcome equivalence or superiority.

The Physical AI Adaption: 21 CFR Part 312 \cite{cfr312-adapt} addresses this
through its Phase 0 simulation validation requirement, which allows Physical
AI systems to demonstrate safety and preliminary efficacy through simulation
before any real patient contact. The 24-hour simulation evidence from the
Clinical Trial Site Complete Documentation \cite{site-docs} provides a model
for how such Phase 0 validation could be conducted at scale.

\subsection{Gaps and Opportunities}
\label{subsec:gaps-opportunities}

Several regulatory gaps and ambiguities exist for Physical AI implementations
that this National Platform addresses. First, the lack of specific FDA guidance
for Physical AI clinical trials means that sponsors must navigate multiple
existing guidances and regulations without a unified reference. This platform
provides that unified reference through its comprehensive body of adapted
standards and documentation.

Second, jurisdictional overlaps between California and federal requirements
create compliance complexity, particularly for data privacy where CMIA, CCPA,
and HIPAA all may apply to the same data. The site establishment documentation
\cite{site-docs} addresses this by integrating both jurisdictions into a single
compliance framework.

Third, classification challenges arise for multi-function Physical AI systems
that simultaneously qualify as medical devices, electronic data systems, and
research tools. The adapted regulatory standards provide classification
guidance, with the 21 CFR Part 312 adaptation \cite{cfr312-adapt} establishing
which Physical AI functions require IND coverage and which fall under device
or operational tool categories.

Fourth, workforce transition considerations must be addressed as Physical AI
enables treating more patients with more robots and fewer workers. California
labor law, federal employment regulations, and healthcare staffing requirements
all intersect with the workforce implications of Physical AI adoption. The
implementation strategy in Section~\ref{sec:implementation} addresses these
workforce considerations through phased deployment with training and role
evolution programs.

\begin{table}[H]
\centering
\caption{Regulatory Gaps and National Platform Solutions}
\label{tab:gaps-solutions}
\small
\begin{tabularx}{\textwidth}{X X}
\toprule
\textbf{Regulatory Gap} & \textbf{National Platform Solution} \\
\midrule
No unified Physical AI trial guidance & Three adapted standards (ICH, CFR 50, CFR 312) \\
Dual-jurisdiction compliance complexity & Integrated site documentation package \\
Multi-function classification challenges & Physical AI definitions across all adaptations \\
Workforce transition regulatory uncertainty & Phased implementation with training programs \\
Absence of Physical AI safety standards & PSL (3 dimensions) and USL (4 dimensions) \\
No multi-site Physical AI infrastructure & MCP servers and federated learning \\
\bottomrule
\end{tabularx}
\end{table}

The regulatory landscape for Physical AI oncology trials is complex but
navigable. By building on existing frameworks rather than creating entirely new
ones, the adapted standards in this platform \cite{ich-adapt, cfr50-adapt,
cfr312-adapt} provide a credible and practical path forward that the
pharmaceutical and regulatory industries can adopt with confidence.
